\documentclass[11pt,a4paper]{article}
\usepackage[margin=27mm]{geometry}
\usepackage[T1]{fontenc}
\usepackage{lmodern}
\usepackage{amsmath,amssymb,amsthm}
\usepackage{hyperref}
\hypersetup{colorlinks=true,linkcolor=blue,urlcolor=blue}
\setlength{\parindent}{0pt}
\setlength{\parskip}{7pt}
\newcommand{\gra}{\operatorname{gra}}
\title{Existence of maximal monotone extensions\\Proof of Conjecture 00000008847}
\author{ziangni-sys}
\date{4 October 2026}
\begin{document}
\maketitle

\textbf{Result.} Both assertions follow from a general maximal-extension
theorem for pairwise conditions on graphs. This theorem permits arbitrary
source and target sets. In particular, it covers monotone operators into
a continuous dual space, as well as monotone operators on inner-product
spaces. No completeness assumption is needed.

\section*{1. General graph formulation}
Let $X,Y$ be arbitrary sets and let $C(p,q)$ be any predicate on pairs of
points $p,q\in X\times Y$. No symmetry, reflexivity, or algebraic
condition on $C$ is assumed. Call $S\subseteq X\times Y$ \emph{admissible}
if
\[
  \forall p,q\in S,\quad C(p,q).
\]
A set-valued map $A:X\to\mathcal P(Y)$ has graph
\[
  \gra A=\{(x,u):u\in A(x)\}.
\]
Call $A$ admissible when its graph is admissible. A map $B$ extends $A$
when $A(x)\subseteq B(x)$ for every $x$, equivalently
$\gra A\subseteq\gra B$. Empty values and empty graphs are allowed.
A map is \emph{maximal admissible} when it is admissible and every
admissible extension of it is equal to it.

\textbf{General theorem.} Every admissible map admits a maximal
admissible extension. Moreover, a map is maximal admissible if and only
if its graph is a maximal element of the family of admissible graphs
ordered by inclusion.

\section*{2. The chain-union lemma}
Let $\mathcal C$ be a family of admissible graphs linearly ordered by
inclusion. Then $U=\bigcup_{S\in\mathcal C}S$ is admissible.

Indeed, choose $p,q\in U$. There are $S,T\in\mathcal C$ with $p\in S$
and $q\in T$. Either $S\subseteq T$ or $T\subseteq S$. Thus the two
points lie together in one admissible graph, which gives $C(p,q)$.
For the empty chain, the union is empty and admissibility is vacuous.
This argument uses only that admissibility is a condition on each pair
of graph points.

\newpage
\section*{3. Existence and graph-order maximality}
Fix an admissible map $A$ and consider
\[
  \mathcal M_A=\{S\subseteq X\times Y:\gra A\subseteq S,
                \ S\text{ is admissible}\},
\]
ordered by inclusion. This family is nonempty because it contains
$\gra A$. A nonempty chain in $\mathcal M_A$ has an upper bound in
$\mathcal M_A$ given by its union: the union contains $\gra A$ and is
admissible by Section 2. The empty chain has $\gra A$ as an upper bound.
Zorn's lemma therefore gives a maximal element $M\in\mathcal M_A$.

Set
\[
  B(x)=\{u\in Y:(x,u)\in M\}.
\]
Then $\gra B=M$, and $B$ is an admissible extension of $A$.
If an admissible map $D$ extends $B$, its graph contains $\gra A$ and
belongs to $\mathcal M_A$. Maximality gives $\gra D=M=\gra B$, so $D=B$.
This proves the extension assertion.

For the characterization, the constructions
\[
  A\longmapsto\gra A,\qquad
  S\longmapsto\bigl[x\mapsto\{u:(x,u)\in S\}\bigr]
\]
are inverse bijections between all set-valued maps $X\to\mathcal P(Y)$
and all subsets of $X\times Y$. They preserve inclusion and
admissibility. Therefore graph-order maximality is exactly the absence
of a proper admissible map extension, in both directions.\hfill$\square$

\section*{4. Standard monotonicity specializations}
\textbf{Continuous-dual operators.} Let $E$ be a real normed space and
$E^*$ its continuous dual, the space of continuous real linear maps
$E\to\mathbb R$. Take $X=E$, $Y=E^*$, and
\[
  C\bigl((x,u),(y,v)\bigr)\quad\Longleftrightarrow\quad
  (u-v)(x-y)\geq0.
\]
The admissible maps are precisely the usual monotone set-valued
operators $A:E\to\mathcal P(E^*)$. The general theorem gives a maximal
monotone extension of every such operator and the exact graph-order
characterization. This applies to every real Banach space; completeness
is unnecessary. Evaluation here is actual evaluation of continuous
linear functionals, not an assumed surrogate pairing.

\textbf{Inner-product operators.} For a real inner-product space $H$,
take $X=Y=H$ and
\[
  C\bigl((x,u),(y,v)\bigr)\quad\Longleftrightarrow\quad
  \langle x-y,u-v\rangle\geq0.
\]
This recovers both conclusions for monotone operators
$H\to\mathcal P(H)$, in particular for Hilbert spaces.
The general theorem also applies to other source/target pairings by
inserting their defining pairwise monotonicity predicate.

\newpage
\section*{5. Scope and formal correspondence}
The English and Chinese conjecture both assert existence of maximal
extensions and maximality under graph inclusion. The general theorem
proves these two assertions for every pairwise graph compatibility
condition, with explicit standard monotone-operator instances in
Section 4. It does not restrict the original claim to equal source and
target spaces or to Hilbert spaces.

No continuity of the set-valued operator, nonempty-domain, full-domain,
or completeness assumption is used. Zorn's lemma uses classical choice.
The result asserts existence, not uniqueness or an effective procedure
for selecting an extension.

The accompanying Lean 4.19.0 project pins Mathlib to commit
\begin{center}\small
\texttt{c44e0c8ee63ca166450922a373c7409c5d26b00b}.
\end{center}
The principal definitions and results in \texttt{lean/Main.lean} are:
\begin{itemize}
\item \texttt{graph}, \texttt{ofGraph}, and both inverse identities for
arbitrary source and target types;
\item \texttt{IsAdmissibleGraph}, the actual condition
$\forall p,q\in S,\ C(p,q)$, and \texttt{admissible\_iff}, its
pointwise map formulation;
\item \texttt{graph\_subset\_iff}, the equivalence between valuewise
map extension and graph inclusion;
\item \texttt{admissible\_sUnion}, the chain-union proof, and
\texttt{exists\_maximal\_graph}, Zorn's lemma with the prescribed graph
as seed;
\item \texttt{maximal\_graph\_iff}, the precise equivalence with absence
of proper admissible map extensions, and
\texttt{exists\_maximal\_extension}, the map existence result.
\end{itemize}

The general final theorem \texttt{conjecture\_00000008847} combines both
conclusions for arbitrary $X$, $Y$, and $C$. No hypothesis asserts an
extension or a maximal element. The source then defines the genuine
inner-product and continuous-dual compatibility predicates and proves
both specialized theorems, \texttt{InnerProduct.result} and
\texttt{Duality.result}. In the dual instance the target is Lean's type
of continuous real linear maps, and \texttt{Duality.monotone\_iff}
identifies the condition with $(u-v)(x-y)\geq0$.

All three final theorems are kernel checked and audited. There are no
additional axioms, incomplete terms, or \texttt{native\_decide}.
The standard logical axioms, including the classical choice used by
Zorn's lemma, are recorded in \texttt{VERIFICATION.md}. No auxiliary
numerical computation is involved.
\end{document}
